\begin{table}[t]
\centering
\caption{Standardized Effect Sizes: Patent Rejection and Inventor Mobility}
\label{tab:sde}
\begin{adjustbox}{max width=\textwidth}
\begin{tabular}{lcccccc}
\hline\hline
Outcome & $\hat{\beta}$ & SE & SD($Y$) & SDE & SE(SDE) & Classification \\
\hline
\textit{Panel A: Pooled} & & & & & & \\
\quad Interstate mobility (full sample) & 0.0110$^{***}$ & (0.0025) & 0.327 & 0.0337 & (0.0076) & Small positive \\
 & & & & & & \\
\textit{Panel B: Heterogeneous} & & & & & & \\
\quad Interstate mobility (solo inventors) & 0.0249$^{**}$ & (0.0098) & 0.394 & 0.0632 & (0.0249) & Moderate positive \\
\quad Interstate mobility (team inventors) & 0.0086$^{***}$ & (0.0025) & 0.322 & 0.0267 & (0.0078) & Small positive \\
\hline\hline
\end{tabular}
\end{adjustbox}
\begin{itemize}[leftmargin=*,noitemsep,topsep=2pt] \item \textit{Notes:} \textbf{Country:} United States. \textbf{Research question:} Does patent application rejection at the USPTO cause inventors to relocate across state lines, redistributing knowledge workers geographically? \textbf{Policy mechanism:} The USPTO assigns patent applications to examiners who vary in grant propensity within technology-specific art units; stricter examiners reject more applications, destroying location-specific intellectual property rents and triggering job search in other states. \textbf{Outcome definition:} Binary indicator equal to one if the inventor's next patent filing (within 10 years) lists a different US state than the current filing. \textbf{Treatment:} Binary --- application rejected (abandoned) versus granted (issued). \textbf{Data:} USPTO PAIR via Google BigQuery, 2002--2014 filings, inventor-application level, 4,452,425 observations. \textbf{Method:} 2SLS using leave-one-out examiner grant rate within art-unit $\times$ filing-year as instrument; clustered standard errors at art-unit $\times$ year level. \textbf{Sample:} US utility patent applications with resolved outcomes (granted or abandoned), examiners with $\geq$2 decisions per art-unit-year, art-unit-years with $\geq$3 examiners. SDE $= \hat{\beta} / \text{SD}(Y)$ where SD($Y$) is the sample standard deviation of the mobility indicator. Classification refers to magnitude, not statistical significance: Large ($|$SDE$| > 0.15$), Moderate ($0.05$--$0.15$), Small ($0.005$--$0.05$), Null ($< 0.005$). $^{***}p<0.01$, $^{**}p<0.05$, $^{*}p<0.1$. \end{itemize}
\end{table}
